\section{How shape becomes mass, CG, and inertia}
For a constant-radius hollow tube, the structural volume is conceptually $\pi(r_o^2-r_i^2)\ell$, multiplied by bulk material density. Its CG lies halfway along its uniform length, and hollow-cylinder inertia formulas distinguish rotation about the long axis from rotation about a transverse axis. Internal rings use the same kind of volume accounting with their own dimensions. A mass object instead supplies explicit mass and an approximate envelope for inertia.

Curved \class{SymmetricComponent} shapes use 128 axial divisions. At each division the code evaluates the external radius at both ends and treats the segment as a conical frustum. A hollow shell subtracts an inner frustum, using wall slope to convert specified thickness into radial material depth; a filled component uses zero inner radii. Segment volumes, first moments, and axial/transverse inertias are accumulated, followed by the parallel-axis shift to the component CG.

The same traversal separately accumulates full external volume, wetted surface area, side-projected planform area, and planform centroid. Full volume describes the outer shape regardless of whether the part is hollow. It is used in aerodynamic CP relationships. Material volume describes what is actually made of the selected material. Substituting one for the other would make a hollow nose's aerodynamic CP incorrectly depend on its wall thickness.

Fin properties are reduced from the outline and supporting root. Material volume includes the profile's thickness factor; mounting tabs and fillets add their own volumes and centroids. Repeated fins, rings of pods, and motor clusters require transformed instance locations before aggregation. Moving a component sideways or rotating an assembly can affect inertia even if total mass and axial CG barely change.

Mass and CG overrides operate at their selected scope. A component-only mass override retains its children as separate contributions. A subtree mass override suppresses the children's aggregate mass contribution; a subtree CG override relocates the children's aggregate center according to the code's override semantics. Unit-inertia approximations remain tied to component geometry. These options are useful for reconciling a measured vehicle, but they are not a procedure for identifying a full measured inertia tensor.

\sourceline{\source{core/src/main/java/info/openrocket/core/rocketcomponent/SymmetricComponent.java}{calculateProperties}; \source{core/src/main/java/info/openrocket/core/rocketcomponent/FinSet.java}{fin mass, tabs, and fillets}; \source{core/src/main/java/info/openrocket/core/masscalc/MassCalculation.java}{structure/motor traversal and overrides}; \source{core/src/main/java/info/openrocket/core/masscalc/MassCalculator.java}{mass calculation entry points}.}

\section{Motor curves, propulsion timing, and database identity}
A thrust-curve motor stores paired time/thrust samples and mass/CG information. \class{ThrustCurveMotor} locates the surrounding sample interval and linearly interpolates thrust and center-of-mass data at the requested motor time. The same object can be queried at several provisional solver times without changing its state. Calculated metadata such as total impulse, average thrust, maximum thrust, and estimated burn duration supports selection and reporting; the designation is not a replacement for the curve.

Loaders handle the different information available in RASP and RockSim motor formats. They validate and normalize input, and may derive mass evolution when the source lacks a full measured mass history. Manufacturer normalization, motor digests, curve sets, and the database finder help distinguish similar names and alternate measurements. A nominal ``same motor'' selection can still refer to different curve data, so reproduction should retain the actual resolved identity.

The installation's \class{IgnitionEvent} can be launch, the lower stage's ejection or burnout, automatic, or never, with an additional configured delay. Automatic selects launch for a launch stage and ejection-triggered ignition for an upper stage. The running \class{MotorClusterState} tracks when that condition actually occurred. A motor's thrust and remaining mass use elapsed time since its own ignition, while the flight branch uses simulation time. Delay/ejection and stage separation are distinct event paths.

Thrust is summed over active motor installations and their physical instances. The model reads supplied thrust curves; it does not solve chamber pressure, nozzle flow, combustion instability, or an electrical ignition circuit. A cluster configuration repeats motor geometry and thrust; it does not imply a modeled ignition-failure distribution unless an explicit extension supplies one.

\sourceline{\source{core/src/main/java/info/openrocket/core/motor/ThrustCurveMotor.java}{ThrustCurveMotor}; \source{core/src/main/java/info/openrocket/core/motor/IgnitionEvent.java}{IgnitionEvent}; \source{core/src/main/java/info/openrocket/core/database/motor/ThrustCurveMotorSet.java}{ThrustCurveMotorSet}; \source{core/src/main/java/info/openrocket/core/simulation/MotorClusterState.java}{MotorClusterState}.}

\section{Atmospheric quantities used by the solver}
The atmosphere supplies pressure and temperature; \class{AtmosphericConditions} derives dry-air density as $\rho=P/(RT)$ with $R=287.053\ \mathrm{J\,kg^{-1}K^{-1}}$. In this implementation the sound speed is the linear approximation $a=165.77+0.606T$ in m/s, with $T$ in kelvin. Kinematic viscosity is $(3.7291\times10^{-6}+4.9944\times10^{-8}T)/\rho$. It is therefore inaccurate to describe this code as evaluating the exact ideal-gas sound-speed expression or full Sutherland viscosity at every sample.

\class{ExtendedISAModel} uses layer boundaries at 0, 11, 20, 32, 47, 51, 71, and 84.852 km, with prescribed base temperatures and hydrostatic pressure relations. Custom launch conditions anchor the lower atmosphere and connect to the higher layer structure. Its constructor restricts a custom first altitude to below 11 km and validates positive absolute temperature and pressure. An interpolation wrapper caches sampled atmospheric values for repeated queries.

The standalone \class{AtmosphericConditions} default is 293.15 K, whereas the standard \class{ExtendedISAModel} sea-level baseline is 288.15 K. Code that constructs a fresh conditions object for a static analysis is not necessarily using the same atmosphere as a flight simulation at its configured launch site. The source also records known approximation limits at high altitude. Humidity, precipitation, and detailed meteorological assimilation are not modeled by these classes.

\section{Wind profiles, gravity, and Earth geometry}
The average wind model combines mean speed/direction with seeded correlated pink-noise turbulence. Turbulence intensity and standard deviation describe variability around the mean; they are not a sequence of independent draws at each RK evaluation. The multilevel model stores altitude-dependent wind levels and interpolates profile behavior, with an explicit AGL/MSL altitude reference. Wind-profile import and editing prepare this model before execution.

Gravity uses a latitude-dependent WGS-style surface magnitude and an inverse-square altitude correction. The geodetic strategy separately controls how local displacement becomes geographic position: \code{FLAT} uses meters-per-degree scaling and zero Coriolis acceleration; \code{SPHERICAL} uses great-circle geometry; \code{WGS84} uses an ellipsoid/Vincenty direct solution. The latter two include the implementation's Coriolis term. Selecting a more detailed map geometry does not replace the rest of the force model.

\sourceline{\source{core/src/main/java/info/openrocket/core/models/atmosphere/AtmosphericConditions.java}{AtmosphericConditions}; \source{core/src/main/java/info/openrocket/core/models/atmosphere/InterpolatingAtmosphericModel.java}{InterpolatingAtmosphericModel}; \source{core/src/main/java/info/openrocket/core/models/wind/PinkNoiseWindModel.java}{PinkNoiseWindModel}; \source{core/src/main/java/info/openrocket/core/models/wind/MultiLevelPinkNoiseWindModel.java}{MultiLevelPinkNoiseWindModel}; \source{core/src/main/java/info/openrocket/core/models/gravity/WGSGravityModel.java}{WGSGravityModel}.}
